\section{Computation in terms of a log resolution}\label{log_resolution}
We use the results of the previous two sections in order to describe Hodge ideals of $\jepPubBold{Q}$-divisors in terms of log resolutions.
Let $X$ be a smooth variety, $h\in\jepPubScript{O}_X(X)$ a nonzero function, $H = {\rm div}(h)$, and $\alpha\in\jepPubBold{Q}_{>0}$. We are interested in computing $I_k(D)$, where $D=\alpha H$.
As always, let $Z={\rm Supp}(D)$ and $\beta = 1 - \alpha$.

Let $f\colon Y\to X$ be a log resolution of the pair $(X, D)$ that is an isomorphism over $U=X\smallsetminus Z$,  
and denote $g=h\circ f\in\jepPubScript{O}_Y(Y)$. 
We fix a positive integer $\ell$ such that $\ell\alpha\in\jepPubBold{Z}$. 
As usual, we consider 
$$p\colon V={\bf Spec}\,\jepPubScript{O}_U[y]/(y^{\ell}-h^{-\ell\alpha})\longrightarrow U$$ 
and the inclusion $j\colon U\hookrightarrow X$. By assumption, we also have an open immersion
$\iota\colon U\hookrightarrow Y$ such that $f\circ\iota=j$. 
 By considering the decompositions of 
$$j_+p_+\jepPubScript{O}_V\simeq f_+\iota_+p_+\jepPubScript{O}_V$$
into isotypical components, we conclude that we have a filtered isomorphism
\begin{equation}\label{birational}
\jepPubScript{M}(h^{-\alpha})\simeq f_+\jepPubScript{M}(g^{-\alpha}).
\end{equation}

We now denote $G=f^*D$, and consider on $Y$ the complex introduced in Section~\ref{complex_for_SNC}:
$$C^{\jepPubBullet}_{g^{-\alpha}}(-\lceil G\rceil):\,\,0\longrightarrow\jepPubScript{O}_Y(-\lceil G\rceil)\otimes_{\jepPubScript{O}_Y} \jepPubScript{D}_Y\longrightarrow\jepPubScript{O}_Y(-\lceil G\rceil)\otimes_{\jepPubScript{O}_Y}\Omega^1_Y(\log E)\otimes_{\jepPubScript{O}_Y}\jepPubScript{D}_Y$$
$$\longrightarrow\cdots\longrightarrow\jepPubScript{O}_Y(-\lceil G\rceil)\otimes_{\jepPubScript{O}_Y}\omega_Y(E)\otimes_{\jepPubScript{O}_Y}\jepPubScript{D}_Y\longrightarrow 0,$$
where $E = (f^*D)_{\rm red}$.
This is placed in degrees $-n,\ldots,0$, and if $x_1,\ldots,x_n$ are local coordinates on $Y$, then its differential is given by
$$\eta\otimes Q\longmapsto d\eta\otimes Q+\sum_{i=1}^n (dx_i\wedge \eta)\otimes\partial_iQ-\big(\alpha\cdot {\rm d\log}(g)\wedge\eta\big)\otimes Q.$$





\begin{theorem}\label{formula_log_resolution}
With the above notation, the following hold:
\begin{enumerate}
\item[(i)] For every $p\neq 0$ and every $k\in \jepPubBold{Z}$, we have
$$R^pf_*\big(C^{\jepPubBullet}_{g^{-\alpha}}(-\lceil G\rceil)\otimes_{\jepPubScript{D}_Y}\jepPubScript{D}_{Y\to X}\big)=0$$
and  
$$R^pf_*F_k\big(C^{\jepPubBullet}_{g^{-\alpha}}(-\lceil G\rceil)
\otimes_{\jepPubScript{D}_Y}\jepPubScript{D}_{Y\to X}\big)=0.$$
\item[(ii)] For every $k\in\jepPubBold{Z}$, the natural inclusion induces an injective map
$$ R^0f_*F_k\big(C^{\jepPubBullet}_{g^{-\alpha}}(-\lceil G\rceil)\otimes_{\jepPubScript{D}_Y}\jepPubScript{D}_{Y\to X}\big)\lhook\joinrel\longrightarrow R^0f_*\big(C^{\jepPubBullet}_{g^{-\alpha}}(-\lceil G\rceil)\otimes_{\jepPubScript{D}_Y}\jepPubScript{D}_{Y\to X}\big).$$
\item[(iii)] We have a canonical isomorphism
$$R^0f_*\big(C^{\jepPubBullet}_{g^{-\alpha}}(-\lceil G\rceil)\otimes_{\jepPubScript{D}_Y}\jepPubScript{D}_{Y\to X}\big)\simeq \jepPubScript{M}_r(h^{-\alpha})$$ 
that induces for every $k\in\jepPubBold{Z}$ an isomorphism
$$R^0f_*F_{k-n}\big(C^{\jepPubBullet}_{g^{-\alpha}}(-\lceil G\rceil)\otimes_{\jepPubScript{D}_Y}\jepPubScript{D}_{Y\to X}\big)\simeq h^{-\alpha}\omega_X\big((k+1)Z-\lceil D\rceil\big)\otimes_{\jepPubScript{O}_X}I_k^\prime (D)$$
$$\simeq h^{\beta}\omega_X\big(kZ + H \big)\otimes_{\jepPubScript{O}_X}I_k(D).$$
\end{enumerate}
\end{theorem}

\begin{proof}
It follows from Lemma~\ref{decomp_filtration}, and from the definition of its filtration, that $\jepPubScript{M}_r(g^{-\alpha})$ is 
a direct summand of a right Hodge $\jepPubScript{D}$-module on $Y$. By Saito's strictness of the filtration of (push-forwards of) such $\jepPubScript{D}$-modules, it follows that for all $k,p\in\jepPubBold{Z}$ the canonical map
$$R^pf_*F_k\big(\jepPubScript{M}_r(g^{-\alpha})\overset{\jepPubBold{L}}{\otimes}_{\jepPubScript{D}_Y}\jepPubScript{D}_{Y\to X}\big)\longrightarrow R^pf_*\big(\jepPubScript{M}_r(g^{-\alpha})\overset{\jepPubBold{L}}{\otimes}_{\jepPubScript{D}_Y}\jepPubScript{D}_{Y\to X}\big)$$
is injective, and its image is equal to
$$F_kR^pf_*\big(\jepPubScript{M}_r(g^{-\alpha})\overset{\jepPubBold{L}}{\otimes}_{\jepPubScript{D}_Y}\jepPubScript{D}_{Y\to X}\big)$$
(see the discussion in Section~\ref{Hodge_modules}).

On the other hand, note that if write 
$G=\alpha\cdot {\rm div}(g)=\sum_i\alpha_iE_i$, then $-\lceil\alpha_i\rceil+\alpha_i\not\in\jepPubBold{Z}_{<0}$ for all $i$.
We may thus apply Proposition~\ref{prop_SNC_complex} for the divisor $G$, with $T=\lceil G\rceil$.
Using Proposition~\ref{Hodge_ideals_SNC} as well,
we see that $C^{\jepPubBullet}_{g^{-\alpha}}(-\lceil G\rceil)$ is filtered quasi-isomorphic to
$\jepPubScript{M}_r(g^{-\alpha})$, hence 
$$R^pf_*\big(C^{\jepPubBullet}_{g^{-\alpha}}(-\lceil G\rceil)\otimes_{\jepPubScript{D}_Y}\jepPubScript{D}_{Y\to X}\big)\simeq R^pf_*\big(\jepPubScript{M}_r(g^{-\alpha})\overset{\jepPubBold{L}}{\otimes}_{\jepPubScript{D}_Y}\jepPubScript{D}_{Y\to X}\big)\quad\text{and}$$
$$R^pf_*F_k\big(C^{\jepPubBullet}_{g^{-\alpha}}(-\lceil G\rceil)\otimes_{\jepPubScript{D}_Y}\jepPubScript{D}_{Y\to X}\big)\simeq 
R^pf_*F_k\big(\jepPubScript{M}_r(g^{-\alpha})\overset{\jepPubBold{L}}{\otimes}_{\jepPubScript{D}_Y}\jepPubScript{D}_{Y\to X}\big).$$

Finally, by the definition of push-forward for right $\jepPubScript{D}$-modules we have
$$R^pf_*\big(\jepPubScript{M}_r(g^{-\alpha})\overset{\jepPubBold{L}}{\otimes}_{\jepPubScript{D}_Y}\jepPubScript{D}_{Y\to X}\big)\simeq 
H^pf_+\jepPubScript{M}_r(g^{-\alpha}),$$
and by (\ref{birational}) this is $0$ if $p\neq 0$, and is canonically isomorphic to $\jepPubScript{M}_r (h^{-\alpha})$ if $p=0$. 
The assertions in the proposition follow by combining all these facts.
\end{proof}
 
\begin{remark}[{\bf Local vanishing}]
The statement in Theorem \ref{formula_log_resolution} i) is a generalization of the Local Vanishing theorem 
for multiplier ideals \cite[Th. 9.4.1]{Lazarsfeld}, in view of the calculation in Proposition 9.1 below.
\end{remark}





